\section*{Capítulo \ref{ch:g9:ions} --- Os íons e as soluções iônicas}

\begin{solution}{exo:g9:ions:1}
Um \omterm{def:g7:atoms-and-molecules:atom}{átomo} ou grupo de \omterm{def:g7:atoms-and-molecules:atom}{átomos} que perdeu ou ganhou \omterm{def:g9:inside-the-atom:nucleus}{elétrons} e tem carga.
Um \omterm{def:g9:ions:ion}{cátion} é positivo (perdeu \omterm{def:g9:inside-the-atom:nucleus}{elétrons}), um \omterm{def:g9:ions:ion}{ânion} é negativo (ganhou
\omterm{def:g9:inside-the-atom:nucleus}{elétrons}).
\end{solution}

\begin{solution}{exo:g9:ions:2}
\ce{Mg^{2+}}: 12 \omterm{def:g9:inside-the-atom:proton}{prótons}, 10 \omterm{def:g9:inside-the-atom:nucleus}{elétrons}.
\end{solution}

\begin{solution}{exo:g9:ions:3}
\ce{KCl}, \ce{MgCl2}, \ce{CaO}.
\end{solution}

\begin{solution}{exo:g9:ions:4}
A água salgada contém \omterm{def:g9:ions:ion}{íons}, \ce{Na+(aq)} e \ce{Cl-(aq)}, que se movem e
transportam carga. As \omterm{def:g7:atoms-and-molecules:molecule}{moléculas} de açúcar não têm carga.
\end{solution}

\begin{solution}{exo:g9:ions:5}
\ce{CuSO4}, \ce{AlCl3}, \ce{Na2CO3}.
\end{solution}

\begin{solution}{exo:g9:ions:6}
Ferro(III), \ce{Fe^{3+}}: \ce{Fe^{3+}(aq) + 3OH-(aq) -> Fe(OH)3(s)}.
\end{solution}

\begin{solution}{exo:g9:ions:7}
Sulfato: 2 \omterm{def:g9:inside-the-atom:nucleus}{elétrons} a mais. Nitrato: 1.
\end{solution}

\begin{solution}{exo:g9:ions:8}
Branco com o nitrato de prata: \ce{Cl-}. Azul com o hidróxido de sódio:
\ce{Cu^{2+}}. Cloreto de cobre(II), \ce{CuCl2}.
\end{solution}

\begin{solution}{exo:g9:ions:9}
Quatro: um de cada lado (à esquerda, à direita, em cima, embaixo).
\end{solution}

\begin{solution}{exo:g9:ions:10}
\ce{Fe^{2+}(aq) + 2OH-(aq) -> Fe(OH)2(s)}.
\end{solution}

\begin{solution}{exo:g9:ions:11}
Cloreto de zinco: os \omterm{def:g9:ions:ion}{íons} sódio não dão \omterm{def:g9:ions:precipitate}{precipitado} com o hidróxido de
sódio, os \omterm{def:g9:ions:ion}{íons} zinco dão um branco:
\ce{Zn^{2+}(aq) + 2OH-(aq) -> Zn(OH)2(s)}.
\end{solution}

\begin{solution}{exo:g9:ions:12}
\ce{CaCl2}: 2000 \omterm{def:g9:ions:ion}{íons} cloreto. Cargas: $1000 \times (+2) + 2000 \times
(-1) = 0$.
\end{solution}

\begin{solution}{pb:g9:ions:1}
\textbf{1.} Ferro(II), \ce{Fe^{2+}}: \ce{Fe(OH)2}.

\textbf{2.} Ferro(III), \ce{Fe^{3+}}: \ce{Fe(OH)3}.

\textbf{3.} O \ce{Fe^{2+}} tem um a mais: $26 - 2 = 24$ \omterm{def:g9:inside-the-atom:nucleus}{elétrons}, contra
$26 - 3 = 23$ do \ce{Fe^{3+}}.

\textbf{4.} \ce{Fe^{3+}(aq) + 3OH-(aq) -> Fe(OH)3(s)}.

\textbf{5.} Ela sai do poço com \omterm{def:g9:ions:ion}{íons} ferro(II); em contato com o \omterm{def:g4:air-a-mixture-of-gases:air}{ar},
eles viram \omterm{def:g9:ions:ion}{íons} ferro(III), que formam o sólido alaranjado.

\textbf{6.} \ce{Fe(OH)3}: $(+3) + 3 \times (-1) = 0$.

\textbf{7.} Não no momento em que ela sai do poço: os \omterm{def:g9:ions:ion}{íons} ferro estão
dissolvidos e passam pelo filtro. Depois que o sólido alaranjado se
formou, sim: ele é um \omterm{def:g9:ions:precipitate}{precipitado}, retido por um filtro.

\textbf{8.} Não: é uma \omterm{def:g6:pure-substances-and-mixtures:mixture}{mistura} (água e \omterm{def:g9:ions:ion}{íons} dissolvidos). Recém-tirada e
límpida, ela é homogênea.

\textbf{9.} $56 \times 1.67 \times 10^{-27} = \qty{9.35e-26}{kg}$.

\textbf{10.} $\qty{0.30}{mg} = \qty{0.30e-6}{kg} = \qty{3.0e-7}{kg}$.

\textbf{11.} Um \omterm{def:g9:ions:ion}{íon} difere do \omterm{def:g7:atoms-and-molecules:atom}{átomo} por dois ou três \omterm{def:g9:inside-the-atom:nucleus}{elétrons}, cuja
massa é desprezível ao lado da massa do \omterm{def:g9:inside-the-atom:nucleus}{núcleo}.

\textbf{12.} $3.0 \times 10^{-7} / 9.35 \times 10^{-26} \approx 3.2
\times 10^{18}$ \omterm{def:g9:ions:ion}{íons} ferro por litro.
\end{solution}
